\section{结论}\label{sec:conclusion}
本文对题目的三个问题给出如下回答。

\textbf{问题一。} 有限库存下不能直接使用算法 1。它不检查库存，售罄商品会持续占用推荐位，既违反零缺货推荐的要求，也浪费了展示其他商品的机会。算例\ref{ex:q1}中，算法 1 相对一个简单可行策略的收入损失在 $T\ge14$ 时超过 30\%，$T=20$ 时为 47.08\%，并随 $T$ 增大趋于 100\%。课程实例上，库存同比例缩小到一半时，算法 1 平均每条路径推荐 \ThetaHalfW{} 次缺货商品。

\textbf{问题二。} 含库存约束的离线模型是流体线性规划\eqref{eq:lp}，它在题目模型中加入期望库存约束，并把推荐决策放松为比例。该模型给出任意在线策略期望收入的上界，课程实例的最优值为 706，影子价格 $\beta^{*}=(9,0,0)$ 表明商品 1 是唯一的稀缺资源，其机会成本来自类型 1 顾客改看商品 3 的代价。

\textbf{问题三。} \LPDual{} 每期重新求解一次剩余销售期的流体线性规划，以库存对偶价格作为机会成本，在有货商品中按调整后收益推荐。策略在每条路径上满足库存约束，每期至多求解一个含 $mn$ 个变量的连续线性规划。五份数据上，\LPDual{} 的平均和最大期望遗憾为 \GavgDual{} 和 \GmaxDual，均低于库存感知贪心（\GavgGreedy、\GmaxGreedy）与有限差分策略（\GavgFD、\GmaxFD）。相对贪心的平均收入提高了 \DualGreedyAvg{}，与预知到达序列最优策略的差距缩小了 \DualSeqGapReductionPct\%。

数值实验还给出以下几点认识。
\begin{enumerate}[label=(\arabic*)]
\item \textbf{购买波动限制了可达收入。} 五份数据上，流体上界 706 与预知顾客顺序时的最优期望收入 \ERSeq{} 之间相差 \FluidGap。商品 1 的随机需求是一个具体来源。这个差额占 \LPDual{} 平均期望遗憾的 \FluidShare\%；从该策略到预知序列最优决策的差额为 \LavgDual。
\item \textbf{同比例放大课程实例时，机会成本的作用更明显。} 在按比例放大的流体模型中，贪心路径与流体最优解的收入差为约 $0.0057T$。在随机购买模拟中，$T=\ScaleTmax$ 时 \LPDual{} 相对贪心的平均配对增收为 $\ScaleGreedyGapMax\pm\ScaleGreedyGapMaxHw$（95\% 置信区间半宽）。最大规模随机配置中，先对每个实例求单期决策时间的第 95 百分位数，再对实例取平均，结果为 \LargeTimeP{} 毫秒（原实验环境）。
\item \textbf{库存条件会改变展示失误的代价。} 在五份课程数据中，逐期损失分解表明，\Greedy{} 的主要失误是向类型 1 过早出售商品 1，\LPFD{} 的主要失误是期末预留过多。沿库存缩放网格观察，算法 1 的缺货展示随库存收紧而增多，其收入损失在中度紧张时达到峰值。
\item \textbf{集中到达会使固定到达率预测失准。} 到达顺序为 $2\to1\to3$ 时，依赖到达率的策略为不会再出现的类型 2 顾客预留库存，期望收入降至 \BlockDual；改用剩余人数预测后恢复到 \BlockCount。
\item \textbf{流体近似看不到需求波动带来的价值。} 流体模型按期望需求预留库存，而课程实例中多留一件的代价约为少留一件的 9 倍，最优预留量远低于期望需求，因此在 $1\to2\to3$ 顺序下，准确的人数预测反而使 \LPCount{} 多留、收入降至 \BlockOTTCount。另一方面，对偶价格使混合使用的推荐集合恰好无差异，库存紧张时（$\theta=0.7$）按即时收益打破并列会多卖；只有一个推荐位时，流体模型又把安全库存视为零价值。这些情形下，LP 策略与随机动态规划之间有 1--7 的差距。
\end{enumerate}

本文的局限与可能的改进方向：购买模型假定顾客不在推荐集合内替代；需求预测只使用固定到达率或总人数；对偶价格在退化点不唯一，评分并列时的规则会影响结果。后续可以考虑在线估计到达率，在并列时按剩余库存排序，或在状态空间较小时用近似动态规划修正流体模型的机会成本。
